\section{Self-Modifying Software：从 Prompt 工程到系统工程}

OpenClaw 最具辨识度的特征是“可自我修改”。这并不是模型突然有了意识，而是系统层允许 agent 在约束内重写自身行为规则。

\subsection{“告诉它不满意什么”与“它去改系统”之间差了什么}

传统对话式 AI 里，用户反馈通常只影响下一轮文本。OpenClaw 的路径是：反馈可触发对配置、文档、脚本、工具链路的变更。

\begin{importantbox}{自修改路径的工程化表达}
访谈可抽象为一个四步机制：
\begin{enumerate}
    \item 用户反馈被转成可执行意图；
    \item agent 定位相关代码或配置；
    \item 通过工具链提交修改并验证；
    \item 变更结果写回记忆和规则文件。
\end{enumerate}
这四步做对，系统才从“会聊”进入“会进化”。
\end{importantbox}

\subsection{为什么这一机制既强大又危险}

\begin{knowledgebox}{强大来源}
\begin{itemize}
    \item 反馈闭环缩短：用户不需等待人工产品迭代；
    \item 局部优化可积累：系统逐步贴合个人工作流；
    \item 复用性增强：修过一次的流程可迁移到相似任务。
\end{itemize}
\end{knowledgebox}

危险同样明显：错误的 feedback 可能被“制度化”进系统，攻击者也可借 prompt injection 诱导写入长期规则。

\subsection{本章小结}
Self-modifying software 的本质是把“学习”从权重层下放到系统层。它提高了演化速度，也把安全和治理压力前置到了产品设计阶段。

